% resumo na língua vernácula (obrigatório)
\setlength{\absparsep}{18pt} % ajusta o espaçamento dos parágrafos do resumo
\begin{resumo}[Resumo]

A análise tradicional da conectividade funcional do cérebro frequentemente colapsa a dimensão temporal, gerando grafos estáticos que mascaram a evolução dinâmica das redes neurais, o chamado "Dinoma" (Dynome). Além disso, a extração de redes resolvidas no tempo a partir de dados de eletrofisiologia de alta densidade (LFP) esbarra no fenômeno físico da condução de volume e no gargalo computacional associado ao processamento estocástico de Big Data. Esta tese de qualificação apresenta o NECTA (Neural Connectivity Time-resolved Analysis), um framework computacional interativo de alto desempenho projetado para a exploração metodológica do mesoconectoma dinâmico. O NECTA emprega uma arquitetura paralela baseada em Dask e armazenamento otimizado nativo de nuvem em Zarr para orquestrar simulações estocásticas em larga escala, fatiamento de tensores multidimensionais e modelos nulos topológicos, superando de maneira linear as limitações críticas de memória RAM (Out-Of-Memory) de fluxos analíticos tradicionais. O framework mitiga ativamente os artefatos de condução volumétrica (zero-phase lag) através de processamentos no sinal analítico de Hilbert, utilizando estimadores complexos baseados em fase, como o Índice de Atraso de Fase Ponderado (wPLI) e a Coerência Imaginária (iCoh). Para inferência causal estrita, implementou-se a Coerência Direcionada Parcial (PDC) em um módulo auto-otimizado dinamicamente pelo Critério de Informação Bayesiana (BIC), controlando penalizações de complexidade em regressões autorregressivas vetoriais (MVAR). O sistema foi validado in silico (contra ruído sintético e condução modelada) e aplicado empiricamente a um dataset de 48 canais do córtex visual primário felino (Área 17, Área 18 e Zona de Transição), registrado sob anestesia geral e paralisia (isolando oscilações intrínsecas de modulações cognitivas atencionais). Os resultados in vivo demonstraram que o córtex visual exibe intensas e ativas flutuações topológicas de alta frequência, atingindo picos transientes de eficiência de "Mundo Pequeno" (Small-Worldness) precisamente coordenados durante a estimulação visual por grades em movimento. A análise direcional na banda Low Gamma (30-59 Hz) estabeleceu a Área 17 como um nó Fonte (Driver) persistente e a Área 18 como um Sumidouro (Sink) estável, enquanto a Zona de Transição atuou como um robusto modulador dinâmico do fluxo inter-hemisférico, alternando sua direcionalidade causal sob estimulação. O NECTA democratiza análises avançadas em neuroinformática, entregando um ecossistema interativo (Dash/Cytoscape) reprodutível focado em Open Science que habilita a tradução ágil e sistêmica de sinais eletrofisiológicos brutos em descobertas conectômicas cronológicas.

\vspace{\onelineskip}
 \noindent
 \textbf{Palavras-chaves}: Conectividade Funcional Dinâmica; Eletrofisiologia; Computação de Alto Desempenho; Coerência Direcionada Parcial; Córtex Visual.
\end{resumo}

% ---
% resumo em inglês
\begin{resumo}[Abstract]
	\begin{otherlanguage*}{english}
	
	Traditional analysis of brain functional connectivity often collapses the temporal dimension, generating static graphs that obscure the dynamic evolution of neural networks, known as the "Dynome". Furthermore, extracting time-resolved networks from high-density electrophysiology (LFP) is heavily hindered by the physical phenomenon of volume conduction and the computational bottleneck associated with stochastic Big Data processing. This qualifying thesis presents NECTA (Neural Connectivity Time-resolved Analysis), an interactive high-performance computational framework designed to explore the dynamic mesoconnectome. NECTA employs a Dask-based parallel architecture and cloud-native Zarr-optimized storage to orchestrate large-scale stochastic simulations, multidimensional tensor slicing, and topological null models, linearly overcoming the critical RAM (Out-Of-Memory) limitations of traditional analytical pipelines. The framework actively mitigates volume conduction artifacts (zero-phase lag) through Hilbert analytical signal processing, utilizing phase-based complex estimators such as the Weighted Phase Lag Index (wPLI) and Imaginary Coherence (iCoh). For strict causal inference, Partial Directed Coherence (PDC) was implemented within a module dynamically auto-optimized by the Bayesian Information Criterion (BIC), controlling complexity penalties in multivariate autoregressive regressions (MVAR). The system was validated in silico (against synthetic noise and modeled conduction) and empirically applied to a 48-channel dataset from the feline primary visual cortex (Area 17, Area 18, and the Transition Zone), recorded under general anesthesia and paralysis (isolating intrinsic oscillations from top-down cognitive modulations). In vivo results demonstrated that the visual cortex exhibits intense and active high-frequency topological fluctuations, reaching transient peaks in Small-Worldness efficiency precisely coordinated during visual stimulation by moving gratings. Directed analysis in the Low Gamma band (30-59 Hz) established Area 17 as a persistent Driver node (Source) and Area 18 as a stable Sink, while the Transition Zone acted as a robust dynamic modulator of interhemispheric flow, alternating its causal directionality upon stimulation. NECTA democratizes advanced neuroinformatics analyses, delivering a reproducible, Open Science-oriented interactive ecosystem (Dash/Cytoscape) that enables the agile and systemic translation of raw electrophysiological signals into chronological connectomic discoveries.
	
	\vspace{\onelineskip}
	\noindent 
	\textbf{Keywords}: Dynamic Functional Connectivity; Electrophysiology; High-Performance Computing; Partial Directed Coherence; Visual Cortex.
	\end{otherlanguage*}
\end{resumo}